% Clean Springer Nature-compatible LaTeX manuscript source.
% Drop-in compatible with Springer Nature `sn-jnl.cls`; structured with standard
% `article` + `natbib`/`booktabs`/`hyperref` packages so any standard TeX setup
% can compile directly without requiring an external proprietary class file.
\documentclass[11pt,a4paper]{article}

\usepackage[T1]{fontenc}
\usepackage[utf8]{inputenc}
\usepackage{lmodern}
\usepackage{geometry}
\geometry{margin=1in}
\usepackage{amsmath,amssymb}
\usepackage{booktabs}
\usepackage{array}
\usepackage{enumitem}
\usepackage[numbers,sort&compress]{natbib}
\usepackage[hidelinks]{hyperref}

\title{Verifiable Multi-Agent Fault-Tolerant Adjudication System with Cryptographic Causal Provenance}

\author{Vaibhav Yadav}

\date{November 2026}

\begin{document}
\maketitle

\begin{abstract}
Multi-agent Large Language Model (LLM) workflows in high-stakes adjudication face four tightly coupled reliability threats: unsafe approvals caused by unconstrained voting over heterogeneous role checks, silent fault propagation under crashes or malformed outputs, stale-certificate replay after source facts mutate, and full-pipeline recomputation cost whenever any input field changes. We present the \emph{Verifiable Multi-Agent Fault-Tolerant Adjudication System with Cryptographic Causal Provenance}, a hybrid architecture that enforces a strict trust boundary between untrusted ML/LLM specialist proposers and a deterministic verification kernel. Under the frozen \texttt{composition-v1} specification, four role-isolated specialists (\texttt{policy}, \texttt{evidence}, \texttt{risk}, \texttt{compliance}) emit structured domain assessments validated against authoritative source bindings and evaluated via a non-overridable mandatory rule gate (\texttt{substantive\_exclusion} and \texttt{strict\_audit}), audited by an independent source-recomputed verifier (\texttt{independent\_verifier}). Every artifact transition is recorded in an Ed25519-signed, SHA-256 content-addressed provenance graph (\texttt{EventDAG}) with an external anti-truncation anchor (\texttt{checkpoint.json}), enabling exact causal lineage queries and certificate-driven selective incremental recovery (\texttt{CertificateSystem}).

On the frozen synthetic held-out benchmark ($N=104$ unique claims across 8 template families; $32$ \texttt{v1}$\rightarrow$\texttt{v2} update paths) and live Groq execution with \texttt{openai/gpt-oss-20b} ($N=36$ stratified claims plus $12$ paired update paths; $325$ API calls, $171{,}888$ tokens), the verified system achieves $1.0000$ exact-match accuracy ($104/104$ offline; $36/36$ live) and $0.0000$ unsafe approvals across clean runs, $1{,}040$ multi-seed fault-injection trials, and live LLM runs. By contrast, non-gated voting baselines commit $70/80$ ($0.8750$) to $73/80$ ($0.9125$) unsafe approvals offline; a strong policy-veto majority baseline still commits $27/80$ ($0.3375$) unsafe approvals on clean claims (McNemar $p = 1.490\times 10^{-8}$) and $406/800$ ($0.5075$) across fault sweeps; and live unverified \texttt{LLM\_AS\_JUDGE} ($30/36$) and \texttt{SINGLE\_AGENT\_E2E\_LLM} ($28/36$) baselines misclassify boundary or mixed-exclusion cases. Cryptographic verification detects $9/9$ tamper, replay, splice, and truncation attacks and blocks $7/7$ stale-certificate replay scenarios, while selective recovery preserves $100.0\%$ decision equivalence on all $32$ offline and $12$ live updates, reducing live LLM calls $4\times$ ($48 \rightarrow 12$). An audited empirical appendix documents a fine-tuned \texttt{DeBERTa-v3-small} advisory dependency classifier (\texttt{Candidate V2}), exposing an \texttt{nli\_context} shortcut confound that motivates keeping learned dependency predictors strictly advisory.
\end{abstract}

\textbf{Keywords:} multi-agent adjudication, cryptographic provenance, selective recovery, mandatory rule gate, Ed25519 audit log, fault tolerance

\section{Introduction}
Large Language Model (LLM) agents are increasingly structured into collaborative multi-agent pipelines for coding, triage, compliance verification, and structured decision workflows~\citep{wu2023autogen,hong2024metagpt,du2024debate}. Yet recent empirical studies demonstrate that multi-agent LLM systems exhibit failure rates ranging from $41\%$ to over $86\%$ due to role drift, broken inter-agent handoffs, unverified assertions, and silent error cascading~\citep{cemri2025mast}. In regulated adjudication domains---such as healthcare insurance claim processing, financial eligibility screening, or safety certification---these coordination failures carry asymmetric risk: approving an ineligible request (an \emph{unsafe approval}) is far costlier than deferring an ambiguous claim to human \texttt{REVIEW}.

We identify four structural gaps in standard multi-agent adjudication designs:
\begin{enumerate}[leftmargin=*,itemsep=2pt]
    \item \textbf{Exclusion-blind voting and category errors.} Treating heterogeneous domain checks (coverage limits, invoice completeness, duplicate screening, identity consent) as interchangeable votes for a binary accept/reject consensus allows non-policy specialists to outvote a substantive policy exclusion, or blurs procedural deferrals (\texttt{REVIEW}) with substantive denials (\texttt{REJECT}).
    \item \textbf{Trusting free-form or judge LLMs without source-bound invariants.} Asking an unconstrained LLM-as-a-judge or single-agent classifier to synthesize evidence leaves the final decision vulnerable to schema truncation, malformed outputs, instruction override, or precedence confusion~\citep{zheng2023judge}.
    \item \textbf{Stale-certificate replay under dynamic updates.} Real claims mutate as missing invoices arrive or claim amounts are amended (\texttt{v1}$\rightarrow$\texttt{v2}). Without cryptographic dependency bindings tying decisions to exact source artifact digests and versions, cached approvals from an earlier epoch can be replayed against stale facts.
    \item \textbf{Expensive full-pipeline restarts vs.\ verification overhead.} Re-executing every specialist on every edit wastes external model calls, yet naive caching cannot prove causal independence unless lineage is explicitly tracked and verified~\citep{mokhov2018build,moreau2013prov}.
\end{enumerate}

To address these gaps, we design and evaluate the \emph{Verifiable Multi-Agent Fault-Tolerant Adjudication System with Cryptographic Causal Provenance}. Its core architectural principle is a hard privilege boundary: \emph{untrusted machine learning and LLM specialists generate typed proposals, whereas deterministic host code owns schema validation, mandatory obligation composition, certificate issuance, cryptographic signing, and incremental recovery}.

\paragraph{Contributions.}
\begin{itemize}[leftmargin=*,itemsep=2pt]
    \item \textbf{Role-Isolated Obligation Composition (\texttt{composition-v1}):} We replace interchangeable voting with disjoint domain contracts across four specialist roles (\texttt{policy}, \texttt{evidence}, \texttt{risk}, \texttt{compliance}), three-valued obligations (\texttt{SAT}, \texttt{FAIL}, \texttt{UNK}), impossible-consequence rejection, and deterministic precedence modes (\texttt{substantive\_exclusion} and \texttt{strict\_audit}) enforced without trust in model-reported confidence scores.
    \item \textbf{Independent Source-Recomputed Certification:} A standalone verifier (\texttt{independent\_verifier}) shares no decision or composition code with the runtime adjudicator and re-derives every role obligation and final decision directly from stored source ancestry before committing decisions.
    \item \textbf{Ed25519 Signed Causal Provenance (\texttt{EventDAG}):} Every mutation, specialist execution, recovery trigger, certificate transition, and committed decision is recorded in a canonical JSON, SHA-256 content-addressed DAG signed by Ed25519 keys with an externally retained \texttt{checkpoint.json} anti-truncation anchor, alongside explicit documentation of unestablished guarantees.
    \item \textbf{Certificate-Driven Selective Recovery:} Upon \texttt{v1}$\rightarrow$\texttt{v2} claim updates, \texttt{CertificateSystem} invalidates only affected subgraph dependencies (\texttt{VALID}$\rightarrow$\texttt{INVALID}$\rightarrow$\texttt{PENDING\_REVERIFICATION}$\rightarrow$\texttt{SUPERSEDED}) and re-executes only invalidated specialist branches while reusing unchanged \texttt{v1} certificates.
    \item \textbf{Frozen Specification, Multi-Arm Evaluation, and Honest Empirical Audit:} We release a hash-locked synthetic benchmark (\texttt{spec-v1}, \texttt{protocol-v1}; $128$ distinct claims, $32$ update paths) and evaluate deterministic rules, strong non-strawman voting baselines, $1{,}040$ multi-seed fault-injection trials, and live Groq \texttt{openai/gpt-oss-20b} runs ($36$ claims, $12$ recovery paths, $325$ API calls), plus an audited appendix analyzing a fine-tuned \texttt{DeBERTa-v3-small} advisory dependency predictor (\texttt{Candidate V2}).
\end{itemize}

\paragraph{Scope and Domain Boundary.}
All domain rules in the adjudication suite (\texttt{coverage-v1}, \texttt{documents-v1}, \texttt{fraud-screen-v1}, \texttt{identity-consent-v1}) operate over a synthetic provisional healthcare claim schema (\texttt{SYNTHETIC\_PROVISIONAL}); no external patient records, private model weights, or unexecuted future variants appear in the release bundle, and ML components hold zero decision or certificate authority.

\section{Related Work}

\subsection{Multi-Agent LLM Orchestration and Failure Taxonomies}
Frameworks such as AutoGen~\citep{wu2023autogen}, MetaGPT~\citep{hong2024metagpt}, and Multiagent Debate~\citep{du2024debate} structure complex LLM workflows through conversational loops, role-specialized standard operating procedures (SOPs), and peer critique. While structured handoffs reduce ad-hoc drift compared to free-form chat, Cemri et al.~\citep{cemri2025mast} systematically analyze over $1{,}600$ execution traces in the MAST taxonomy and identify $14$ failure modes spanning system specification errors, inter-agent misalignment, and missing or flawed task verification. In contrast to purely prompt-driven role assignment, our architecture enforces disjoint role input scoping at runtime and backs specialist outputs with hard deterministic obligation composition and independent ancestry re-verification.

\subsection{LLM-as-a-Judge and Calibration Limits}
Using LLMs as automated adjudicators or evaluators has gained wide adoption~\citep{zheng2023judge}, accompanied by documented position, verbosity, self-enhancement, and finite-step reasoning biases. When deployed to resolve conflicting role findings or adjudicate numerical policy thresholds (such as exact INR coverage limit equality), an unconstrained LLM judge risks conflating procedural uncertainty with substantive rejection or failing on strict JSON formatting under token/syntax errors. Our live evaluation empirically compares an explicit \texttt{LLM\_AS\_JUDGE} arm and a joint \texttt{SINGLE\_AGENT\_E2E\_LLM} baseline directly against gated role-isolated composition under identical \texttt{openai/gpt-oss-20b} calls.

\subsection{Rule-Gated Verification, Proof-Carrying Code, and Selective Prediction}
Proof-Carrying Code (PCC)~\citep{necula1997pcc} separates an untrusted producer that proposes artifacts and candidate proofs from a small, trusted verifier that independently checks safety invariants before execution. Selective classification~\citep{geifman2017selective} formalizes the safety-coverage trade-off via a learned or calibrated abstention option that lowers risk on accepted predictions at the expense of reduced automated coverage. We unify both ideas for multi-agent adjudication: specialist proposals are untrusted, confidence weights carry no commitment authority, and unvalidated or conflicting states abstain deterministically into human \texttt{REVIEW} rather than risking an unsafe \texttt{APPROVE}.

\subsection{Incremental Build Systems and Selective Recomputation}
Incremental build systems model task dependencies as directed acyclic graphs so that edits recompute only transitive dependents while skipping verified unchanged targets~\citep{mokhov2018build}. Constructive traces and content-addressed keys ensure that rebuilding a clean node produces identical results to a fresh from-scratch execution. Our \texttt{CertificateSystem} applies these invariants to multi-agent AI adjudication, coupling fine-grained fact versioning with explicit certificate invalidation so that unaffected specialist LLM calls are skipped without altering final decisions.

\subsection{Cryptographic Causal Provenance and Tamper-Evident Logging}
W3C PROV-DM~\citep{moreau2013prov} formalizes provenance graphs over entities, activities, and agents (\texttt{used}, \texttt{wasGeneratedBy}, \texttt{wasDerivedFrom}), while tamper-evident logging systems combine cryptographic digests and Merkle accumulators to detect unauthorized edits~\citep{crosby2009tamper}. Certificate Transparency (RFC~6962)~\citep{laurie2013rfc6962} establishes how externally retained Signed Tree Heads (STHs) prevent log equivocation and silent tail truncation, whereas Practical Byzantine Fault Tolerance (PBFT)~\citep{castro1999pbft} achieves state-machine safety across replicated distrusted nodes via quorum agreement. High-speed deterministic Edwards-curve signatures (Ed25519)~\citep{bernstein2012ed25519,josefsson2017rfc8032} provide compact $64$-byte authentication with deterministic nonces. In contrast to distributed BFT consensus~\citep{castro1999pbft}, our system assumes a single-writer local host verifier that bounds faulty specialist proposals deterministically and pairs Ed25519 hash-chained lineage with externally retained \texttt{checkpoint.json} anchors~\citep{crosby2009tamper,laurie2013rfc6962}, while clearly separating operational audit edges from data-dependency lineage during selective recovery.

\section{Threat Model and Security Properties}

\subsection{Adversarial and Fault Capabilities}
We consider a hybrid environment in which LLM specialist backends, transport links, and persisted execution logs may experience both benign operational failures and active tampering:
\begin{itemize}[leftmargin=*,itemsep=2pt]
    \item \textbf{Agent and Transport Faults (\texttt{F01}--\texttt{F05}):} Specialists may exhaust bounded retries due to API timeouts or rate limits (\texttt{crash\_timeout}), emit unparseable prose, invalid JSON, out-of-range confidences, or forbidden decision tokens (\texttt{malformed\_output}), bind outputs to outdated input hashes (\texttt{stale\_binding}), return schema-valid JSON whose decision or copied facts violate domain rules or whose free-text rationale contradicts its own structured fields (\texttt{valid\_but\_wrong}), or act as a single compromised specialist trying to force \texttt{APPROVE} despite an active policy exclusion (\texttt{adversarial\_specialist}).
    \item \textbf{Artifact Replay Attacks:} A caller attempts to commit a previously valid round certificate (\texttt{round-certificate@1}) or re-authorize a round after an upstream source fact mutates (either flipping or preserving the label).
    \item \textbf{Provenance Store Attacks:} An attacker with write access to persisted JSONL logs attempts payload mutation, forged SHA-256 recomputation without valid signing keys, rogue-key re-signing, parent-link manipulation, interior deletion, event reordering, cross-run event splicing, prefix truncation, or rollback against an external checkpoint.
\end{itemize}

\subsection{Established Security Properties and Explicit Trust Boundaries}
Table~\ref{tab:trust-boundary} summarizes both the security guarantees established by the deterministic kernel and the explicit non-claims audited in \texttt{provenance\_verifier.py} and \texttt{fault\_matrix.json}.

\begin{table}[t]
\centering
\caption{Security Guarantees Established vs.\ Explicit Trust-Boundary Limits}
\label{tab:trust-boundary}
\small
\begin{tabular}{@{}p{0.46\linewidth}p{0.48\linewidth}@{}}
\toprule
\textbf{Established by Design} & \textbf{Not Established (Explicit Limits)} \\
\midrule
Tamper evidence on any byte change to payload, parents, sequence, \texttt{run\_id}, or \texttt{artifact\_ref} via SHA-256 digests and Ed25519 signatures. & Real-world identity binding beyond possession of the trusted Ed25519 private key (no external PKI/CA). \\
Cross-run splice rejection via monotonic sequence checks and bound \texttt{run\_id}. & Ground-truth factuality of initial source inputs (``garbage-in'' fact records still sign cleanly). \\
Stale-certificate commit and derivation refusal on mutated \texttt{content\_hash} or superseded versions. & Semantic faithfulness of non-empty free-text \texttt{reasoning\_summary} when \texttt{decision}, \texttt{evidence}, and \texttt{rule\_references} match (\texttt{F04b}). \\
Tail truncation and checkpoint rollback detection when an independently retained \texttt{checkpoint.json} is supplied~\citep{crosby2009tamper,laurie2013rfc6962}. & Tail truncation detection \emph{without} a retained external checkpoint (\texttt{C10} limit: a prefix remains internally valid). \\
Data-dependency lineage parity across selective recovery by separating data links from \texttt{recovery\_started} edges. & Distributed or Byzantine quorum consensus~\citep{castro1999pbft} (single-writer local log; crash between \texttt{events.jsonl} and \texttt{checkpoint.json} leaves a lagging checkpoint gap). \\
\bottomrule
\end{tabular}
\end{table}

\section{System Architecture and Formal Specification}

\subsection{Disjoint Specialist Contracts and Three-Valued Obligations}
Each claim $x$ is partitioned into four disjoint typed fact slices so that no specialist sees facts belonging to another role or any ground-truth oracle label:
\begin{itemize}[leftmargin=*,itemsep=2pt]
    \item \textbf{\texttt{policy} (\texttt{coverage-v1}, obligation \texttt{policy\_eligibility}):} Inspects positive integers \texttt{amount\_inr}, \texttt{coverage\_limit\_inr} and booleans \texttt{hospitalization\_covered}, \texttt{waiting\_period\_satisfied}. Evaluates \texttt{SAT} (\texttt{APPROVE}) iff covered $\land$ waiting\_satisfied $\land$ ($\texttt{amount\_inr} \le \texttt{coverage\_limit\_inr}$), and \texttt{FAIL} (\texttt{REJECT}) otherwise. \texttt{REVIEW} is an illegal role consequence.
    \item \textbf{\texttt{evidence} (\texttt{documents-v1}, \texttt{evidence\_integrity}):} Evaluates \texttt{SAT} (\texttt{APPROVE}) iff \texttt{documents\_available} $\land$ \texttt{invoice\_matches}, else \texttt{FAIL} (\texttt{REVIEW}). \texttt{REJECT} is illegal.
    \item \textbf{\texttt{risk} (\texttt{fraud-screen-v1}, \texttt{screening\_integrity}):} Evaluates \texttt{SAT} (\texttt{APPROVE}) iff $\neg\texttt{duplicate\_claim} \land \neg\texttt{provider\_flagged}$, else \texttt{FAIL} (\texttt{REVIEW}). \texttt{REJECT} is illegal.
    \item \textbf{\texttt{compliance} (\texttt{identity-consent-v1}, \texttt{regulatory\_controls}):} Evaluates \texttt{SAT} (\texttt{APPROVE}) iff \texttt{identity\_verified} $\land$ \texttt{consent\_available}$, else \texttt{FAIL} (\texttt{REVIEW}). \texttt{REJECT} is illegal.
\end{itemize}

The host validator (\texttt{assess\_specialist}) maps each role response into $\mathcal{O}_r \in \{\texttt{SAT}, \texttt{FAIL}, \texttt{UNK}\}$. Any execution failure, non-JSON prose, schema violation (e.g., non-finite confidence or coerced decision token), \texttt{IMPOSSIBLE\_CONSEQUENCE}, \texttt{stale\_or\_unbound\_input} binding mismatch, copied \texttt{evidence} mismatch, \texttt{rule\_references} mismatch, or \texttt{decision\_conflicts\_with\_rule} error forces $\mathcal{O}_r = \texttt{UNK}$ (consequence forced to \texttt{REVIEW}), preventing crashes from impersonating either \texttt{APPROVE} or substantive \texttt{REJECT}.

\subsection{Mandatory Rule Composition and Independent Certification}
Given validated obligations $(\mathcal{O}_p, \mathcal{O}_e, \mathcal{O}_r, \mathcal{O}_c)$, \texttt{compose\_specialist\_assessments} evaluates rules strictly top-to-bottom under \texttt{substantive\_exclusion}:
\begin{equation}
\hat{y} =
\begin{cases}
\texttt{REVIEW}  & \text{if } \mathcal{O}_p = \texttt{UNK} \quad (\texttt{POLICY\_UNCERTAIN}),\\
\texttt{REJECT}  & \text{else if } \mathcal{O}_p = \texttt{FAIL} \quad (\texttt{SUBSTANTIVE\_EXCLUSION}),\\
\texttt{REVIEW}  & \text{else if } \exists x \in \{e,r,c\}:\ \mathcal{O}_x \in \{\texttt{FAIL}, \texttt{UNK}\} \quad (\texttt{PROCEDURAL\_DEFERRAL}),\\
\texttt{APPROVE} & \text{if } \mathcal{O}_p = \mathcal{O}_e = \mathcal{O}_r = \mathcal{O}_c = \texttt{SAT} \quad (\texttt{ALL\_OBLIGATIONS\_SATISFIED}).
\end{cases}
\end{equation}
In the stricter \texttt{strict\_audit} mode, any $\mathcal{O}_x = \texttt{UNK}$ for $x \in \{p,e,r,c\}$ halts to \texttt{REVIEW} before evaluating exclusions. Crucially, the certificate verifier \texttt{verify\_composition\_independently} imports \emph{zero} code from \texttt{composition.py}, \texttt{consensus.py}, \texttt{domain.rule\_decision}, or \texttt{cert\_recovery.oracle}: it reads the authoritative source facts and declared bindings directly from DAG ancestry and rejects any commit whose derived decision, gate, bindings, or \texttt{EventDAG} ledger linkage disagree.

\subsection{Ed25519 Signed EventDAG and Certificate Selective Recovery}
The core engine (\texttt{CertificateSystem}) stores immutable JSON snapshots keyed by $(\texttt{artifact\_id}, \texttt{version})$ and SHA-256 content hashes. Every state transition appends an event record (\texttt{signed-event-dag-v1}) binding \texttt{run\_id}, monotonic \texttt{sequence}, \texttt{payload\_digest}, \texttt{parent\_hashes}, \texttt{signer\_key\_id}, and an Ed25519 signature over canonical domain-separated bytes (\texttt{cert-recovery.event.v1}), persisted to \texttt{events.jsonl} alongside an atomically replaced \texttt{checkpoint.json} (\texttt{cert-recovery.checkpoint.v1}).

When an input fact slice changes from \texttt{v1} to \texttt{v2}, \texttt{change\_dependency} traverses downstream consumers, invalidating affected certificates (\texttt{VALID}$\rightarrow$\texttt{INVALID}), planning targeted recovery (\texttt{PENDING\_REVERIFICATION}), re-executing \emph{only} invalidated specialist roles, reusing unchanged \texttt{v1} role outputs, issuing \texttt{round-certificate@2}, and superseding the prior certificate (\texttt{SUPERSEDED}). Recovery ordering edges (\texttt{recovery\_started}) link historical finality into the audit log; causal data-lineage queries explicitly distinguish operational recovery ordering edges from data-dependency parent edges so provenance audits remain exact.

\section{Methodology and Frozen Experimental Protocol}

\subsection{Frozen Synthetic Dataset and Leakage Isolation (\texttt{spec-v1})}
All primary evaluations run against the hash-locked synthetic benchmark frozen under \texttt{datasets/frozen\_v1/} and \texttt{results/phase5/frozen\_evaluation\_manifest.json}:
\begin{itemize}[leftmargin=*,itemsep=2pt]
    \item \textbf{Split Separation:} $128$ distinct claims partitioned with strict zero overlap across template families, fact hashes, and non-approve boolean signatures into a $24$-claim \texttt{dev} set ($4$ families; $6$ \texttt{APPROVE}, $10$ \texttt{REJECT}, $8$ \texttt{REVIEW}) and a $104$-claim held-out \texttt{test} set ($8$ families; $24$ \texttt{APPROVE}, $46$ \texttt{REJECT}, $34$ \texttt{REVIEW}; $80$ total non-approve cases, $58$ non-reject cases). Notice the test split includes $1$-INR boundary stress cases (\texttt{amount\_inr == coverage\_limit\_inr} vs.\ \texttt{coverage\_limit\_inr + 1}) and mixed policy-exclusion plus procedural-failure cases.
    \item \textbf{Frozen Update Paths:} $32$ \texttt{v1}$\rightarrow$\texttt{v2} claim transitions spanning four strata: \texttt{narrow\_single\_role} ($12$), \texttt{broad\_multi\_role} ($8$), \texttt{no\_change\_edit} ($6$), and \texttt{same\_label\_value\_edit} ($6$).
    \item \textbf{Integrity Preflight:} Every Phase~6 execution script verifies the SHA-256 digests of all $8$ specification, protocol, dev, test, and update files prior to evaluation, aborting on any mismatch. (An external blind-teammate file \texttt{blind\_external\_template.json} remains unpopulated with $0$ submissions and is explicitly reported as \texttt{NOT RUN}.)
\end{itemize}

\subsection{Evaluation Arms, Non-Strawman Baselines, and Statistical Protocol}
To avoid strawman comparisons, the historical exclusion-blind prototype (\texttt{Arm E} in \texttt{legacy} mode) is deprecated as a historical reference only, and primary evaluation tests seven protocol-defined arms plus five ablations (\texttt{protocol-v1}):
\begin{enumerate}[leftmargin=*,itemsep=2pt]
    \item \textbf{\texttt{FULL\_DET} and \texttt{FULL\_LIVE}:} Complete \texttt{composition-v1} system with disjoint role isolation, host validation, mandatory rule gate, independent verifier, and Ed25519 \texttt{EventDAG}, backed either by deterministic role rules (\texttt{FULL\_DET}) or isolated live LLM calls (\texttt{FULL\_LIVE}).
    \item \textbf{\texttt{MAJORITY\_VOTE} and \texttt{CONFIDENCE\_WEIGHTED\_VOTE}:} Unweighted majority (ties to \texttt{REVIEW}) and confidence-weighted voting over specialist outputs without the mandatory exclusion gate.
    \item \textbf{\texttt{MAJORITY\_POLICY\_VETO} (Strong Baseline):} Grants \texttt{policy} mandatory \texttt{REJECT} veto power over eligible exclusions while aggregating non-exclusion disagreements by unweighted majority voting without typed obligation composition.
    \item \textbf{\texttt{SINGLE\_AGENT\_E2E\_LLM} and \texttt{LLM\_AS\_JUDGE} (Live LLM Baselines):} A single monolithic LLM prompt seeing all $10$ claim fields jointly (\texttt{SINGLE\_AGENT\_E2E\_LLM}) and a separate LLM judge adjudicating the four specialist outputs without deterministic gate enforcement (\texttt{LLM\_AS\_JUDGE}).
\end{enumerate}

We report overall exact-match accuracy, unsafe-approval rate ($\text{UnsafeApprovalRate} = \#\{\hat{y}_i=\texttt{APPROVE} \land y_i\neq\texttt{APPROVE}\} / \#\{y_i\neq\texttt{APPROVE}\}$), false-reject rate, coverage ($\#\{\hat{y}_i\in\{\texttt{APPROVE},\texttt{REJECT}\}\}/N$), review abstention rate, and selective accuracy~\citep{geifman2017selective}. Proportions carry $95\%$ Wilson score intervals~\citep{wilson1927probable}; exact paired binary comparisons use two-sided exact McNemar tests~\citep{mcnemar1947note} alongside family-clustered paired bootstraps ($B=2{,}000$). Multi-seed offline sweeps run across $10$ fixed seeds ($\{101,202,\dots,1010\}$) over the $5$ fault classes and are reported as exploratory.

\subsection{Live Execution Calibration (\texttt{openai/gpt-oss-20b} via Groq)}
Because the hosted Groq endpoint enforces an $8{,}000$ tokens-per-minute (TPM) free-tier ceiling, Phase~4 calibrated outbound pacing at $\texttt{min\_call\_interval\_sec}=5.0$, \texttt{temperature}$=0.0$, bounded retries ($\le 2$ on transient \texttt{429}/\texttt{5xx}/timeouts, fail-fast on \texttt{400}/\texttt{401}/\texttt{403}/\texttt{404}), and rate-hint backoff. Under this ceiling, Phase~6 executes the frozen Option-B stratified live subset ($N=36$ balanced claims across all $8$ test families: $8$ \texttt{APPROVE}, $16$ \texttt{REJECT}, $12$ \texttt{REVIEW}; plus $12$ stratified update paths: $3$ per update kind).

\section{Experimental Evaluation}

\subsection{Phase 4 Live Pilot and Fault Containment}
In the Phase~4 live development pilot (\texttt{results/phase4/pilot\_summary.json}), \texttt{groq:openai/gpt-oss-20b} completed $24$ clean specialist calls across $6$ dev claims ($1$ \texttt{APPROVE}, $2$ \texttt{REJECT}, $3$ \texttt{REVIEW}), matching the deterministic oracle on $6/6$ claims with $0$ \texttt{UNK} coercions. Combining clean calls and live-wrapped injected faults consumed $58$ total calls ($18{,}641$ prompt $+ 8{,}188$ completion $= 26{,}829$ total tokens, $\$0.00$ cost). Across $8$ controlled injected fault scenarios (\texttt{results/phase4/fault\_matrix.json}), $8/8$ matched expected containment behavior: \texttt{F01\_crash\_timeout\_exhausted}, \texttt{F02a\_malformed\_non\_json}, \texttt{F02b\_schema\_invalid\_decision\_token}, \texttt{F03\_stale\_binding}, and \texttt{F04a\_valid\_json\_wrong\_decision} were each mapped to \texttt{UNK}/\texttt{REVIEW}; \texttt{F05a} and \texttt{F05b} (adversarial compromised \texttt{evidence} under policy exclusion) resolved to \texttt{REJECT} under \texttt{substantive\_exclusion} and \texttt{REVIEW} under \texttt{strict\_audit}; and \texttt{F04b\_valid\_json\_unfaithful\_reasoning} behaved honestly as an explicit accepted limit when structured fields matched authoritative rules despite contradictory prose.

\subsection{Primary Held-Out Offline Evaluation ($N=104$ Clean and $1{,}040$ Multi-Seed Faults)}
Table~\ref{tab:offline-main} summarizes held-out test performance on clean claims (seed $101$) and across the exploratory $10$-seed fault-injection sweep ($104 \times 10 = 1{,}040$ trials, $800$ non-approve instances).

\begin{table}[t]
\centering
\caption{Frozen Held-Out Offline Evaluation ($N=104$ Clean, Seed $101$; and $1{,}040$ Exploratory Fault Trials)}
\label{tab:offline-main}
\small
\begin{tabular}{@{}lcccccc@{}}
\toprule
\textbf{Arm} & \textbf{Exact Acc.} & \textbf{95\% Wilson CI} & \textbf{Unsafe App.} & \textbf{Coverage} & \textbf{Sel.\ Acc.} & \textbf{McNemar $p$ (Acc / Safe)} \\
\midrule
\multicolumn{7}{@{}l}{\textit{Panel A: Clean Held-Out Test Set ($N=104$, $80$ non-approve)}} \\
\texttt{FULL\_DET} & \textbf{1.0000} & [0.9644, 1.0000] & \textbf{0.0000} ($0/80$) & 0.6731 & 1.0000 & Baseline ref \\
\texttt{MAJORITY\_POLICY\_VETO} & 0.7404 & [0.6486, 0.8150] & 0.3375 ($27/80$) & 0.9327 & 0.7216 & $1.490\times 10^{-8}$ / $1.490\times 10^{-8}$ \\
\texttt{MAJORITY\_VOTE} & 0.2981 & [0.2186, 0.3919] & 0.8750 ($70/80$) & 0.9038 & 0.2553 & $2.118\times 10^{-22}$ / $1.694\times 10^{-21}$ \\
\texttt{CONFIDENCE\_WEIGHTED\_VOTE} & 0.2692 & [0.1933, 0.3616] & 0.9125 ($73/80$) & 0.9327 & 0.2474 & $2.647\times 10^{-23}$ / $2.118\times 10^{-22}$ \\
\midrule
\multicolumn{7}{@{}l}{\textit{Panel B: Exploratory $10$-Seed Fault Sweep ($N=1{,}040$ trials, $800$ non-approve)}} \\
\texttt{FULL\_DET} & \textbf{0.7048} & [0.6764, 0.7317] & \textbf{0.0000} ($0/800$) & 0.3779 & 1.0000 & Exploratory pooled \\
\texttt{MAJORITY\_POLICY\_VETO} & 0.6019 & --- & 0.5075 ($406/800$) & 0.9471 & 0.5878 & Exploratory pooled \\
\texttt{MAJORITY\_VOTE} & 0.2837 & --- & 0.8375 ($670/800$) & 0.8827 & 0.2702 & Exploratory pooled \\
\texttt{CONFIDENCE\_WEIGHTED\_VOTE} & 0.3288 & --- & 0.8525 ($682/800$) & 0.9510 & 0.3104 & Exploratory pooled \\
\bottomrule
\end{tabular}
\end{table}

On clean claims, \texttt{FULL\_DET} achieves $104/104$ exact matches and $0/80$ unsafe approvals ($0.0000$, $95\%$ Wilson upper bound $0.0458$) with zero false rejections ($0/58$) and $1.0000$ selective accuracy on its $70$ automated \texttt{APPROVE}/\texttt{REJECT} decisions. Even when \texttt{MAJORITY\_POLICY\_VETO} enforces policy rejections, unweighted non-policy majority voting still approves $27/80$ ($0.3375$) procedural-\texttt{REVIEW} claims whenever two non-policy agents pass while one fails ($p=1.490\times 10^{-8}$ vs.\ \texttt{FULL\_DET}; family-clustered bootstrap $95\%$ accuracy difference $[0.0000, 0.5870]$).

Table~\ref{tab:per-fault} disaggregates Seed-$101$ performance across all five fault classes ($N=104$ per class). \texttt{FULL\_DET} sustains $0.0000$ unsafe approvals ($0/80$) and $1.0000$ selective accuracy across every fault regime by routing corrupted inputs cleanly to \texttt{REVIEW} (\texttt{0.6250}--\texttt{0.9038} overall accuracy), whereas all voting baselines exhibit $0.4500$--\texttt{0.9625} unsafe approval rates.

\begin{table}[t]
\centering
\caption{Per-Fault Breakdown on Frozen Test Set ($N=104$ per Fault Class, Seed $101$)}
\label{tab:per-fault}
\small
\begin{tabular}{@{}lcccc@{}}
\toprule
\textbf{Fault Class} & \textbf{\texttt{FULL\_DET}} & \textbf{\texttt{MAJ.\_POLICY\_VETO}} & \textbf{\texttt{MAJORITY\_VOTE}} & \textbf{\texttt{CONF.\_WEIGHTED}} \\
 & (Acc / Unsafe App.) & (Acc / Unsafe App.) & (Acc / Unsafe App.) & (Acc / Unsafe App.) \\
\midrule
\texttt{adversarial\_specialist} & 0.9038 / \textbf{0.0000} & 0.6154 / 0.5000 & 0.2692 / 0.9250 & 0.2404 / 0.9625 \\
\texttt{crash\_timeout}           & 0.6538 / \textbf{0.0000} & 0.6058 / 0.5000 & 0.2788 / 0.8500 & 0.3365 / 0.8500 \\
\texttt{malformed\_output}        & 0.6538 / \textbf{0.0000} & 0.5962 / 0.5250 & 0.2692 / 0.8625 & 0.3365 / 0.8625 \\
\texttt{stale\_binding}           & 0.6250 / \textbf{0.0000} & 0.5962 / 0.5125 & 0.2981 / 0.8750 & 0.2692 / 0.9125 \\
\texttt{valid\_but\_wrong}        & 0.7115 / \textbf{0.0000} & 0.6442 / 0.4500 & 0.3269 / 0.4875 & 0.6058 / 0.5000 \\
\bottomrule
\end{tabular}
\end{table}

\subsection{Live Multi-Arm Evaluation with Groq \texttt{openai/gpt-oss-20b} ($N=36$)}
Across the $36$-claim live Option-B subset and $12$ live recovery paths, the live run consumed $325$ API calls ($116{,}922$ prompt $+ 54{,}966$ completion $= 171{,}888$ total tokens, $\$0.00$ free-tier USD spend; Table~\ref{tab:live-arms}).

\begin{table}[t]
\centering
\caption{Live Option-B Multi-Arm Evaluation on Groq \texttt{openai/gpt-oss-20b} ($N=36$ Stratified Claims)}
\label{tab:live-arms}
\small
\begin{tabular}{@{}lcccccc@{}}
\toprule
\textbf{Live Arm} & \textbf{Exact Acc.} & \textbf{Unsafe App.} & \textbf{Coverage} & \textbf{Sel.\ Acc.} & \textbf{Calls} & \textbf{Total Tokens} \\
\midrule
\texttt{FULL\_LIVE}            & \textbf{1.0000} ($36/36$) & 0.0000 & \textbf{0.6667} ($24/36$) & 1.0000 & 144 & 74{,}886 \\
\texttt{LLM\_AS\_JUDGE}        & 0.8333 ($30/36$)          & 0.0000 & 0.5000 ($18/36$)          & 1.0000 & 36\textsuperscript{a} & 23{,}524\textsuperscript{a} \\
\texttt{SINGLE\_AGENT\_E2E\_LLM} & 0.7778 ($28/36$)        & 0.0000 & 0.4444 ($16/36$)          & 1.0000 & 36  & 17{,}117 \\
\bottomrule
\end{tabular}

\vspace{2pt}
\footnotesize{\textsuperscript{a}\texttt{LLM\_AS\_JUDGE} reused the $144$ validated specialist outputs from \texttt{FULL\_LIVE} and issued $36$ judge calls.}
\end{table}

Inspecting \texttt{results/phase6/raw/live\_claims.jsonl} and \texttt{live\_judge.jsonl} reveals concrete mechanistic causes for the accuracy gaps even at zero unsafe approvals:
\begin{itemize}[leftmargin=*,itemsep=2pt]
    \item \texttt{FULL\_LIVE} achieved $36/36$ ($1.0000$) exact match and $24/36$ ($0.6667$) automated coverage, committing $0$ unsafe approvals and $0$ false rejections.
    \item \texttt{LLM\_AS\_JUDGE} dropped to $30/36$ ($0.8333$) across two failure channels: ($i$) $3$ judge calls failed strict JSON parsing due to unterminated string truncation (`\texttt{ParseError: Unterminated string}`) and fell back safely to \texttt{REVIEW} on true \texttt{REJECT} claims; and ($ii$) on $3$ claims containing confirmed policy exclusions (\texttt{p5\_test\_fam\_test\_mixed\_exclusion\_procedural\_002}, \texttt{...\_006}, and \texttt{...\_single\_hosp\_excl\_003}), the LLM judge returned \texttt{REVIEW} instead of \texttt{REJECT}, violating \texttt{substantive\_exclusion} precedence.
    \item \texttt{SINGLE\_AGENT\_E2E\_LLM} dropped to $28/36$ ($0.7778$) because it returned over-conservative \texttt{REVIEW} on all $8/8$ valid \texttt{APPROVE} boundary/micro claims (\texttt{p5\_test\_fam\_test\_clean\_boundary\_and\_micro\_*}), cutting automated coverage from $0.6667$ to $0.4444$.
\end{itemize}

\subsection{Cryptographic Provenance, Replay Defense, and Selective Recovery}
\paragraph{Phase 2 Stale-Certificate Replay ($7/7$ PASS).}
In \texttt{results/phase2/stale\_certificate\_replay\_matrix.json}, clean issuance and post-recovery commit controls (\texttt{1\_clean\_positive\_control}, \texttt{5\_unrelated\_edit\_control}, \texttt{6\_post\_recovery\_positive\_control}) succeeded, while all $4$ stale replay scenarios (\texttt{2\_stale\_commit\_replay\_label\_changes}, \texttt{3\_stale\_authorization\_replay\_label\_changes}, \texttt{4a\_same\_label\_stale\_replay\_approve}, \texttt{4b\_same\_label\_stale\_replay\_reject}) were refused at commit or authorization even when a numeric fact edit preserved the final boolean outcome.

\paragraph{Phase 3 Ed25519 Provenance Tamper Matrix ($11/11$ PASS).}
On a $41$-event persisted log (\texttt{run\_id=tamper-exp-v1}, \texttt{results/phase3/provenance\_tamper\_matrix.json}), provenance verification accepted the clean reload (\texttt{C01}), detected $9/9$ attacks ($0$ accepted: \texttt{C02} payload byte tamper, \texttt{C03} forged SHA-256 with original signature, \texttt{C04} rogue-key re-signing, \texttt{C05} parent-link alteration, \texttt{C06} interior deletion, \texttt{C07} adjacent reorder, \texttt{C08} cross-run splice, \texttt{C09} tail truncation against retained checkpoint, and \texttt{C11} checkpoint rollback), and recorded the honest \texttt{C10} control (truncation without an external checkpoint remains locally valid). Across $3$ causal cross-check queries (\texttt{facts:evidence@1}, \texttt{results/phase3/causal\_crosscheck.json}), stopping traversal at operational \texttt{recovery\_started} edges yielded $3/3$ exact lineage matches against the artifact DAG.

\paragraph{Phase 6 Selective Recovery and Ablations.}
Table~\ref{tab:recovery} and Table~\ref{tab:ablations} summarize recovery efficiency across the $32$ offline update paths and the five targeted Phase~6 ablations.

\begin{table}[t]
\centering
\caption{Selective Recovery vs.\ Full Rerun ($32$ Offline Updates and $12$ Live Groq Updates)}
\label{tab:recovery}
\small
\begin{tabular}{@{}lcccccc@{}}
\toprule
\textbf{Update Kind} & \textbf{Count} & \textbf{Spec.\ Execs} & \textbf{Nodes Rebuilt} & \textbf{Sel.\ vs.\ Full} & \textbf{Equiv.} & \textbf{Stale (Cert /} \\
 & & \textbf{(Sel.\ / Full)} & \textbf{/ Skipped} & \textbf{Time (ms)} & & \textbf{No-Cert)} \\
\midrule
\multicolumn{7}{@{}l}{\textit{Panel A: $32$ Frozen Offline Update Paths (\texttt{FULL\_DET})}} \\
\texttt{narrow\_single\_role}  & 12 & 1.00 / 4.00 & 4.00 / 3.00 & 23.210 vs.\ 21.490 & 100.0\% & $0 / 12$ \\
\texttt{broad\_multi\_role}    &  8 & 2.12 / 4.00 & 5.12 / 1.88 & 26.766 vs.\ 21.073 & 100.0\% & $0 / 8$  \\
\texttt{no\_change\_edit}      &  6 & 0.00 / 4.00 & 0.00 / 7.00 &  0.721 vs.\ 21.366 & 100.0\% & $0 / 0$  \\
\texttt{same\_label\_value\_edit} & 6 & 1.00 / 4.00 & 4.00 / 3.00 & 23.879 vs.\ 21.140 & 100.0\% & $0 / 6$ \\
\midrule
\textbf{Total / Aggregate}   & \textbf{32} & \textbf{35 / 128} & --- & \textbf{20.008 vs.\ 21.297} & \textbf{100.0\%} & \textbf{0 / 26} \\
\midrule
\multicolumn{7}{@{}l}{\textit{Panel B: $12$ Paired Live Groq \texttt{openai/gpt-oss-20b} Update Paths ($3$ per Kind)}} \\
\texttt{LIVE\_RECOVERY\_PAIRED} & 12 & \textbf{12 / 48} ($4\times$) & --- & API-paced ($5.0\text{s}$/call) & \textbf{100.0\%} & --- \\
\bottomrule
\end{tabular}
\end{table}

Crucially, Table~\ref{tab:recovery} illustrates a realistic cost crossover: when specialist evaluation consists of microsecond local Python rules, verification ($14.59$--$16.86$\,ms) and Ed25519 signing ($8.33$--$10.57$\,ms) exceed deterministic rule compute, making narrow single-role selective recovery marginally slower in CPU wall clock ($23.210$\,ms vs.\ $21.490$\,ms) while no-change edits drop to $0.721$\,ms. Conversely, in live LLM deployment where each call incurs multi-second network latency and token billing, selective recovery preserves $12/12$ ($100.0\%$) decision equivalence against both cached deterministic replay and fresh $4$-agent live reruns while cutting external API calls by $4\times$ ($48 \rightarrow 12$). Meanwhile, removing certificate staleness checks (\texttt{no\_certificates}) causes stale decisions on $26/32$ ($81.25\%$) updates, compared to $0/32$ ($0.0\%$) with active certificates.

\begin{table}[t]
\centering
\caption{Summary of Five Targeted Phase~6 Ablations and Historical Phase~1 Baseline Parity}
\label{tab:ablations}
\small
\begin{tabular}{@{}lp{0.28\linewidth}cc@{}}
\toprule
\textbf{Ablation ID} & \textbf{Isolated Property} & \textbf{With Component} & \textbf{Without Component} \\
\midrule
\texttt{no\_gate} & Unsafe-approval rate (seed $101$ fault) & \textbf{0.0000} ($0/80$) & 0.8875 ($71/80$) \\
\texttt{no\_validation} & Faulted-input containment ($N=104$) & \textbf{0.8654} ($90/104$) & 0.4231 ($44/104$) \\
\texttt{no\_certificates} & Stale-decision rate ($32$ updates) & \textbf{0.0000} ($0/32$) & 0.8125 ($26/32$) \\
\texttt{no\_provenance} & Tamper detection / lineage query & \textbf{9/9}, Supported & 0/9, Unsupported \\
\texttt{no\_recovery} & Specialist calls ($32$ offline updates) & \textbf{35} ($100\%$ equiv.) & 128 ($3.66\times$ calls) \\
\bottomrule
\end{tabular}
\end{table}

All Phase~1 deterministic seed tables also regenerate identically via \texttt{run\_all.py} across $N=150$ synthetic claims ($13/13$ regenerated metrics matching historical artifacts, including \texttt{COMP} $84.0\%$ acc / $38/45$ rejects vs.\ legacy validated voting $72.0\%$ acc / $3/45$ rejects, and factorial ablation $0$ unsafe approvals under gated validation vs.\ $107$--$109$ ungated).

\section{Discussion}
Three design lessons emerge from the joint offline and live evaluations. First, \emph{architectural role isolation plus deterministic precedence beats scaling a single prompt or judge}: on identical \texttt{openai/gpt-oss-20b} inference, splitting claims into disjoint structured role contracts raised exact-match accuracy from $0.7778$ (single-agent) and $0.8333$ (LLM judge) to $1.0000$, eliminating both boundary over-caution and mixed-exclusion errors without any model fine-tuning. Second, \emph{certificate staleness is orthogonal to label stability}: $6/6$ same-label value edits still invalidated source hashes and were re-certified rather than silently reusing \texttt{v1} certificates, whereas $6/6$ no-change edits bypassed specialist execution completely. Third, \emph{provenance verification costs are negligible relative to external model inference}, even though microsecond deterministic baselines highlight those fixed cryptographic costs explicitly.

\section{Limitations and Future Work}
We highlight six explicit boundaries of the current study:
\begin{enumerate}[leftmargin=*,itemsep=2pt]
    \item \textbf{Synthetic Domain Schema and Blind External Split:} All contracts (\texttt{coverage-v1}, \texttt{documents-v1}, \texttt{fraud-screen-v1}, \texttt{identity-consent-v1}) use deterministic synthetic predicates; the external blind-teammate file (\texttt{datasets/blind\_external\_template.json}) contains $0$ submissions (\texttt{NOT RUN}). Generalization to open-ended legal or clinical policies requires formalizing richer statutory predicates.
    \item \textbf{Groq Free-Tier Subset Size:} Under an $8{,}000$ TPM ceiling, live evaluation covered a stratified $36$-claim subset and $12$ update paths rather than the full $104$-claim, $32$-update matrix.
    \item \textbf{Unverified Free-Text Reasoning Summaries (\texttt{F04b}):} While host validation enforces exact decision, copied-fact, and rule-citation parity, free-text \texttt{reasoning\_summary} strings are checked only for non-emptiness; verifying natural-language faithfulness remains open.
    \item \textbf{No-Checkpoint Truncation Limit (\texttt{C10}) and Crash Window:} Detecting tail truncation requires an independently retained \texttt{checkpoint.json}; an atomic two-file write crash between \texttt{events.jsonl} and \texttt{checkpoint.json} leaves a recoverable lagging-checkpoint gap.
    \item \textbf{Single-Writer Trust Assumption:} Ed25519 signing establishes possession of a single writer private key rather than distributed PKI identity or Byzantine quorum consensus.
    \item \textbf{Learned Dependency Predictors Remain Non-Authoritative:} As audited in Appendix~\ref{app:v2}, fine-tuned cross-encoders can still exploit semantic shortcut cues, forbidding their use for authoritative DAG edge creation.
\end{enumerate}

\section{Conclusion}
We presented a verifiable multi-agent adjudication architecture uniting disjoint specialist role composition, independent ancestry-based certification, Ed25519 hash-chained causal provenance, and selective incremental recovery. Across $104$ frozen held-out claims, $1{,}040$ multi-seed fault trials, $32$ update paths, and $36+12$ live Groq \texttt{openai/gpt-oss-20b} executions, hard deterministic enforcement achieved zero unsafe approvals and $100\%$ decision fidelity while cutting live specialist calls $4\times$ during incremental updates and catching $9/9$ provenance attacks and $7/7$ stale replay attempts. By isolating untrusted neural proposals behind cryptographic bindings and source-recomputed rule gates, multi-agent systems gain verifiable safety without abandoning incremental efficiency.

\appendix
\section{Empirical Audit of Advisory Dependency Classifier (\texttt{Candidate V2})}
\label{app:v2}
Beyond runtime adjudication, the repository includes an offline empirical investigation of whether fine-tuned neural classifiers can replace explicit declared dependency declarations in support graphs, trained against \texttt{cross-encoder/nli-deberta-v3-small}~\citep{he2023debertav3} and compared with pretrained MultiNLI-style NLI baselines~\citep{williams2018multinli}. All metrics below come strictly from verified frozen bundle artifacts (\texttt{results/phase7/v2\_appendix\_summary.json}); no fresh training or unexecuted follow-up model runs are claimed.

\begin{table}[t]
\centering
\caption{Audited Performance of \texttt{Candidate V2} vs.\ Pretrained Baseline and Smoke Checks}
\label{tab:v2-audit}
\small
\begin{tabular}{@{}lcccccc@{}}
\toprule
\textbf{Evaluation Split} & \textbf{Model} & \textbf{Accuracy} & \textbf{95\% Wilson CI} & \textbf{Macro-F1} & \textbf{FP / FN} & \textbf{Audit Note} \\
\midrule
Validation ($n=64$) & \texttt{Candidate V2} ($\tau=0.20$) & 0.8750 ($56/64$) & [0.7723, 0.9353] & 0.8730 & 8 / 0 & Selected $\tau=0.20$ \\
Test ($n=96$)       & Original NLI                      & 0.5000 ($48/96$) & ---              & 0.3333 & 48 / 0 & $96/96$ all-positive \\
Test ($n=96$)       & \texttt{Candidate V2}             & 0.9167 ($88/96$) & [0.8441, 0.9572] & 0.9161 & 8 / 0 & All 8 FPs on \texttt{nli\_context} \\
Stress ($n=384$)    & Original NLI                      & 0.8464           & ---              & 0.8426 & 59 / 0 & \texttt{nli\_trap}: $15/64$ FPs \\
Stress ($n=384$)    & \texttt{Candidate V2}             & 0.8281           & ---              & 0.8229 & 66 / 0 & \texttt{nli\_trap}: $64/64$ FPs \\
\bottomrule
\end{tabular}
\end{table}

As Table~\ref{tab:v2-audit} shows, while \texttt{Candidate V2} eliminated the original all-positive failure mode on the $96$-pair test split ($50.00\% \rightarrow 91.67\%$ accuracy, $0$ false negatives), stress testing on $n=384$ rows regressed overall accuracy ($84.64\% \rightarrow 82.81\%$, with $42$ item improvements offset by $49$ regressions) because every single false positive on test ($8/8$) and all $64/64$ \texttt{nli\_trap} stress negatives occurred when the hypothesis semantically restated an observed fact absent declared support edges. Independent smoke verification (\texttt{baseline\_tests.json}) likewise confirmed $5/6$ accuracy for both \texttt{nli-deberta-v3-small} and \texttt{nli-MiniLM2-L6-H768}, failing contrast pair \texttt{nli-003} at $0.9933$ and $0.9957$ contradiction scores, alongside $50.0\%$ lexical identifier overlap across templates. Consequently, \texttt{semantic\_advice.py} enforces strict read-only containment: advisory scores capture explicit $(\texttt{artifact\_id}, \texttt{version}, \texttt{sha256})$ bindings and reject stale replies (\texttt{REJECTED\_STALE}), holding zero authority to create DAG edges or issue certificates.

\bibliographystyle{plainnat}
\bibliography{references}

\end{document}
